% Supplement to the frozen 2026-11-06 v3 manuscript.
% This supplies local calculations, not a replacement residual/table artifact.
\paragraph{An explicit check of the enlarged anchor.}
The prime $3$ is inert and unramified in $K=\mathbb Q(\sqrt2)$.
The displayed point $P$ reduces to $(\sqrt2,0)$, a nonzero point of
order two. Exact doubling gives
\[
 z(2P)=\frac{3396356159088-4135204317756\sqrt2}{3287811625151},
 \qquad v_3(z(2P))=1.
\]
In the integral formal group, multiplication by a $3$-adic unit
preserves positive integral depth, and multiplication by $3$ increases
it by one. If $[a]P$ is formal and nonzero, then $a$ is even and
\[
 v_3(f([a]P))=2v_3(z([a]P))=2\bigl(1+v_3(a/2)\bigr).
\]
The anchor with $C=3^{2^{26}}$ forces
$v_3(f([a]P))\ge2^{26}$. Thus
$0\ne a\in2\,3^{2^{25}-1}\mathbb Z$.
This proves the size estimate used in the precision-$52$ argument,
conditional only on the stated integer-label description of the
allowed points and the anchor's already proved semilocal equivalence.

\paragraph{The support change under anchor translation.}
Put $A=a+b-c$. The square of
$-3A^2+16u^2+96v^2-32w^2+32s^2$ has $27$ monomials involving
one of its three private coordinates. The square of
\[
 -3A^2+U^2-6UV+6UW+15V^2-18VW+7W^2+8sU
\]
has $64$ such monomials. Among the old monomials not involving
private coordinates, the translation removes the constant and the
six degree-two monomials in $a,b,c$, but leaves the degree-four
part unchanged. The stated parent-table overlap data says that the
constant and $a^2,b^2,c^2$ remain present in the other residual
squares, while the three mixed degree-two monomials do not.
The support therefore increases by $64-27-3=34$.
Consequently the parent count $28987$, if verified from its table,
gives $29021$ after translation. This calculation does not replace
the verification of the parent table or its overlap data.
